% `tex.web` §1025: the output routine gets a nest level of its own.
%
%     push_nest; mode:=-vmode; prev_depth:=ignore_depth; mode_line:=-line;
%
% Negative, so it is *internal* vertical mode: `\ifinner` is true there, and
% §211's `print_mode` calls it "internal vertical mode" — which §367's command
% trace consults to decide whether to name the mode at all, since it names it
% only when it has changed since the last line.
%
% This engine runs `\output` on the outer vertical list with `\prevdepth` and
% `\prevgraf` saved rather than opening a box for it, so the mode looked outer
% from the inside: `\ifinner` was false, and the first command *after* the
% routine did not name its mode because as far as the trace could tell it had
% not changed. In trip.log that is `{vertical mode: \looseness}` against our
% bare `{\looseness}`, three hundred lines after the routine that caused it.
\catcode`\{=1 \catcode`\}=2
\tracingonline=1
\tracingcommands=1
\vsize=30pt \maxdepth=2pt
\tracingoutput=0

\output={\message{[inner=\ifinner Y\else N\fi]}\shipout\box255}
\message{[outer=\ifinner Y\else N\fi]}
\hbox{}\vfill\penalty-10000

% The command after the routine names its mode again, because the routine's
% was different.
\count1=1
\hbox{}
\end
